\section[Análisis de trade-offs]{Análisis de trade-offs\\(N-Capas vs.\ Clean/Onion/Hexagonal)}

\subsection{Enunciado}
Un equipo de desarrollo está construyendo un sistema de gestión de inventario. Inicialmente,
eligen una Arquitectura N-Capas tradicional por su simplicidad y familiaridad. Después de un
año, el sistema se ha vuelto un monolito difícil de mantener, donde los cambios en la base de
datos afectan la lógica de negocio y la interfaz de usuario, y las pruebas unitarias son lentas
y frágiles.

Con base en los principios de la Arquitectura Limpia y la Arquitectura Hexagonal, explique
detalladamente qué error fundamental de acoplamiento se está cometiendo y cómo la introducción
de la Regla de Dependencia (Martin) o el uso de Puertos y Adaptadores (Cockburn) resolverían
estructuralmente este problema. ¿Cuál sería el impacto en la testabilidad y la capacidad de
cambiar la base de datos de SQL Server a MongoDB sin afectar la lógica de negocio?

\subsection{Respuesta}

\subsubsection{¿Qué error fundamental de acoplamiento se está cometiendo y cómo la introducción de la Regla de Dependencia (Martin) o el uso de Puertos y Adaptadores (Cockburn) resolverían estructuralmente este problema?}

En una arquitectura N-Capas \emph{clásica} (figura~\ref{fig:n-capas}) hay una separación física
en proyectos, módulos o ensamblados .

\begin{figure}[H]
  \centering
  \begin{tikzpicture}[node distance=8mm]
    \node[figbox, fill=figUiFill, minimum width=12.2cm] (ui) {%
      \textbf{Presentación}\\[1pt]
      {\scriptsize\color{figMuted} Proyecto.Presentacion · vistas y controladores}};
    \node[figbox, fill=figNegFill, minimum width=12.2cm, below=of ui] (neg) {%
      \textbf{Lógica de negocio}\\[1pt]
      {\scriptsize\color{figMuted} Proyecto.Negocio · servicios de aplicación}};
    \node[figbox, fill=figInfraFill, draw=figInfra, minimum width=12.2cm, below=of neg] (dat) {%
      \textbf{Acceso a datos}\\[1pt]
      {\scriptsize\color{figMuted} Proyecto.Datos · ORM, SQL y procedimientos almacenados}};
    \draw[figarr] (ui.south) -- (neg.north);
    \draw[figarr] (neg.south) -- (dat.north);
    \node[font=\scriptsize, text=figMuted, anchor=west]
      at ($(ui.south)!0.5!(neg.north)+(0.28,0)$) {depende de};
    \node[font=\scriptsize, text=figMuted, anchor=west]
      at ($(neg.south)!0.5!(dat.north)+(0.28,0)$) {depende de};
  \end{tikzpicture}
  \caption{Arquitectura N-Capas clásica. Cada capa puede vivir en un ensamblado distinto, pero la dependencia de compilación apunta hacia el detalle inferior.}
  \label{fig:n-capas}
\end{figure}

Esa separación no garantiza el aislamiento conceptual del dominio.
Cuando las capas superiores dependen directamente de detalles concretos (ORM, consultas SQL,
DTOs de persistencia, controles de interfaz), el \textbf{dominio queda acoplado a la
infraestructura} y a los mecanismos de entrega (figura~\ref{fig:acoplamiento})
\parencite{martin2017clean,torre2010guia}.

\begin{figure}[H]
  \centering
  \begin{tikzpicture}[node distance=7mm and 6mm]
    \node[figbox, fill=figNucleoFill, draw=figNucleo, line width=0.9pt,
          minimum width=4.7cm, minimum height=1.7cm] (dom) {%
      \textbf{Dominio}\\[1pt]
      {\scriptsize reglas de inventario}};
    \node[figbox, fill=figBadFill, draw=figBad, minimum width=4.15cm, above=of dom] (ui) {%
      \textbf{Controles de interfaz}\\[1pt]
      {\scriptsize vistas, grids, controladores}};
    \node[figbox, fill=figBadFill, draw=figBad, minimum width=4.15cm, left=of dom] (sql) {%
      \textbf{SqlConnection y ORM}\\[1pt]
      {\scriptsize consultas y esquema}};
    \node[figbox, fill=figBadFill, draw=figBad, minimum width=4.15cm, right=of dom] (dto) {%
      \textbf{DTOs de persistencia}\\[1pt]
      {\scriptsize filas, tablas, documentos}};
    \node[figbox, fill=figBadFill, draw=figBad, minimum width=4.15cm, below=of dom] (sp) {%
      \textbf{Procedimientos almacenados}\\[1pt]
      {\scriptsize reglas escritas en SQL}};
    \draw[figarrbad] (dom) -- (ui);
    \draw[figarrbad] (dom) -- (sql);
    \draw[figarrbad] (dom) -- (dto);
    \draw[figarrbad] (dom) -- (sp);
  \end{tikzpicture}
  \caption{Error de acoplamiento. Las flechas salen del dominio: la lógica de inventario depende de la interfaz y de la persistencia, así que un cambio de esquema atraviesa el negocio y la UI. Elaboración propia.}
  \label{fig:acoplamiento}
\end{figure}

El síntoma del enunciado (cambios en la base de datos que arrastran reglas de negocio y UI;
pruebas lentas y frágiles) suele deberse a una violación combinada de:

\begin{itemize}
  \item \textbf{Dependencias invertidas incorrectamente:} la lógica de negocio importa tipos del
    acceso a datos o del framework web, en lugar de que la infraestructura implemente contratos
    definidos por el dominio \parencite{martin2017clean}.
  \item \textbf{Modelo anémico y lógica dispersa:} las reglas de inventario viven en servicios de
    aplicación o en procedimientos almacenados, mientras las entidades son meros contenedores de
    datos, lo que difumina el límite entre ``negocio'' y ``persistencia''.
  \item \textbf{Acoplamiento transversal:} transacciones, validaciones y mapeos repetidos en cada
    capa generan un monolito lógico en el que ningún cambio local es verdaderamente local.
\end{itemize}

En términos de atributos de calidad, el sistema sacrifica \emph{mantenibilidad} y
\emph{testabilidad} por una simplicidad inicial de despliegue en un solo proceso
\parencite{richards2020fundamentals}.

La Arquitectura Limpia organiza el software en anillos concéntricos: entidades y casos de uso en el
centro; interfaces, persistencia y frameworks en el exterior \parencite{martin2017clean}. La
\textbf{Regla de Dependencia} exige que todo vector de dependencia (importaciones, herencia,
referencias) apunte \emph{hacia adentro}. Nada en el núcleo puede conocer SQL Server, MongoDB,
HTTP ni la UI (figura~\ref{fig:anillos}).

\begin{figure}[H]
  \centering
  \begin{tikzpicture}[font=\footnotesize]
    \def\rA{1.62}
    \def\rB{2.68}
    \def\rC{3.74}
    \def\rD{4.80}
    \fill[figRingFw] (0,0) circle (\rD);
    \fill[figRingAd] (0,0) circle (\rC);
    \fill[figRingUc] (0,0) circle (\rB);
    \fill[figRingEn] (0,0) circle (\rA);
    \draw[figInk, line width=0.55pt] (0,0) circle (\rD);
    \draw[figInk, line width=0.55pt] (0,0) circle (\rC);
    \draw[figInk, line width=0.55pt] (0,0) circle (\rB);
    \draw[figInk, line width=0.55pt] (0,0) circle (\rA);
    \node[align=center, text=white, font=\scriptsize] at (0,0.02) {%
      \textbf{\footnotesize Entidades}\\[2pt]
      Producto\\[-1pt] Almacén\\[-1pt] Movimiento};
    \node at (90:2.15) {Casos de uso};
    \node at (90:3.21) {Adaptadores};
    \node at (90:4.27) {Frameworks};
    \draw[-{Stealth[length=2.5mm,width=1.9mm]}, line width=0.8pt, draw=figNucleo]
      (0,-5.28) -- (0,-1.78);
    \node[anchor=east, font=\scriptsize, text=figNucleo]
      at (-0.22,-5.35) {Regla de dependencia};
    \node[anchor=north west, align=left, font=\scriptsize] at (5.15,3.05) {%
      \FigSwatch{figRingFw}\hspace{1.3mm}\textbf{Frameworks y drivers}\\
      \hspace*{5.1mm}UI, HTTP, SQL Server, MongoDB\\[6pt]
      \FigSwatch{figRingAd}\hspace{1.3mm}\textbf{Adaptadores de interfaz}\\
      \hspace*{5.1mm}controladores y gateways\\[6pt]
      \FigSwatch{figRingUc}\hspace{1.3mm}\textbf{Casos de uso}\\
      \hspace*{5.1mm}\texttt{RegistrarEntradaDeStock}\\[6pt]
      \FigSwatch{figRingEn}\hspace{1.3mm}\textbf{Entidades}\\
      \hspace*{5.1mm}Producto, Almacén, Movimiento%
    };
  \end{tikzpicture}
  \caption{Regla de dependencia. El código solo puede depender hacia el centro: los casos de uso y las entidades no conocen SQL Server, MongoDB, HTTP ni la interfaz. Elaboración propia a partir de \parencite{martin2017clean}.}
  \label{fig:anillos}
\end{figure}

Aplicado al inventario, por ejemplo, el caso de uso \emph{RegistrarEntradaDeStock} dependería de una abstracción
(\texttt{IRepositorioDeProducto} o un puerto equivalente), no de \texttt{SqlConnection} ni de
documentos BSON. La capa de datos implementa esa abstracción y traduce entre el modelo de dominio y
el almacén concreto. Así, un cambio de esquema relacional a documentos no reescribe reglas de
negocio: solo sustituye o añade un adaptador de salida (figura~\ref{fig:inventario}).

\begin{figure}[H]
  \centering
  \begin{tikzpicture}[node distance=11mm]
    \node[figbox, fill=figNucleoFill, draw=figNucleo, line width=0.9pt, minimum width=8.8cm] (uc) {%
      \textbf{RegistrarEntradaDeStock}\\[1pt]
      {\scriptsize caso de uso · reglas de inventario}};
    \node[figbox, fill=white, draw=figNucleo, line width=0.9pt, minimum width=8.8cm, below=of uc] (port) {%
      \texttt{IRepositorioDeProducto}\\[1pt]
      {\scriptsize puerto de salida, definido en el núcleo}};
    \node[figbox, fill=figAdaptFill, draw=figAdapt, text width=4.85cm,
          anchor=north] (sql) at ([xshift=-3.2cm,yshift=-12mm]port.south) {%
      {\scriptsize\texttt{SqlServerInventarioAdapter}}\\[2pt]
      {\scriptsize tablas y transacciones ACID}\\[1pt]
      {\scriptsize\color{figMuted} SQL Server}};
    \node[figbox, fill=figAdaptFill, draw=figAdapt, text width=4.85cm,
          anchor=north] (mongo) at ([xshift=3.2cm,yshift=-12mm]port.south) {%
      {\scriptsize\texttt{MongoInventarioAdapter}}\\[2pt]
      {\scriptsize colecciones y mapeo a documentos}\\[1pt]
      {\scriptsize\color{figMuted} MongoDB}};
    \draw[figarr] (uc.south) -- (port.north);
    \node[font=\scriptsize, text=figMuted, anchor=west]
      at ($(uc.south)!0.5!(port.north)+(0.22,0)$) {depende de};
    \draw[figimpl] (sql.north) -- (port.south -| sql.north);
    \draw[figimpl] (mongo.north) -- (port.south -| mongo.north);
    \node[font=\scriptsize, text=figAdapt] at ($(port.south)+(0,-0.48)$) {implementan};
    \draw[decorate, decoration={brace, amplitude=4.5pt, mirror}, draw=figOk, line width=0.6pt]
      (sql.south west) -- (mongo.south east)
      node[midway, below=8pt, font=\scriptsize, align=center, text=figOk]
      {intercambiables: el caso de uso no se recompila};
  \end{tikzpicture}
  \caption{El caso de uso depende del puerto, no del motor. SQL Server y MongoDB quedan detrás de adaptadores que implementan la misma abstracción.}
  \label{fig:inventario}
\end{figure}

La Arquitectura Hexagonal expresa el mismo principio con la metáfora de \textbf{puertos}
(interfaces que el dominio expone o requiere) y \textbf{adaptadores} (implementaciones concretas
para UI, API, base de datos, mensajería) \parencite{cockburn2005hexagonal}. El núcleo de la
aplicación permanece aislado; los adaptadores son intercambiables.

\begin{itemize}
  \item \textbf{Puerto de entrada:} por ejemplo, \texttt{GestionarInventario} invocado desde una
    API REST o una consola; ambos son adaptadores de entrada distintos sobre el mismo puerto.
  \item \textbf{Puerto de salida:} por ejemplo, \texttt{PersistirMovimientoDeInventario}; un
    adaptador habla con SQL Server y otro con MongoDB sin que el dominio distinga el motor.
\end{itemize}

Estructuralmente, esto rompe el acoplamiento en cadena ``UI $\rightarrow$ servicio
$\rightarrow$ repositorio concreto'' y lo reemplaza por ``adaptador $\rightarrow$ aplicación
$\rightarrow$ puerto'', con las dependencias cerradas hacia el centro (figura~\ref{fig:hexagonal}).

\begin{figure}[H]
  \centering
  \begin{tikzpicture}[font=\footnotesize]
    \def\R{2.15}
    \draw[fill=figNucleoFill, draw=figNucleo, line width=0.9pt]
      (90:\R) -- (150:\R) -- (210:\R) -- (270:\R) -- (330:\R) -- (30:\R) -- cycle;
    \node[align=center] at (0,0) {%
      \textbf{Núcleo}\\[1pt]
      {\scriptsize casos de uso y entidades}};
    \node[figbox, fill=white, draw=figNucleo, line width=0.8pt,
          minimum width=2.55cm, anchor=east] (pin) at (-2.22,0) {%
      {\scriptsize\texttt{Gestionar}}\\
      {\scriptsize\texttt{Inventario}}\\
      {\scriptsize\color{figMuted}entrada}};
    \node[figbox, fill=white, draw=figNucleo, line width=0.8pt,
          minimum width=2.55cm, anchor=west] (pout) at (2.22,0) {%
      {\scriptsize\texttt{Persistir}}\\
      {\scriptsize\texttt{Movimiento}}\\
      {\scriptsize\color{figMuted}salida}};
    \node[figbox, fill=figAdaptFill, draw=figAdapt, minimum width=2.7cm, anchor=east]
      (api) at (-5.15,0.92) {API REST};
    \node[figbox, fill=figAdaptFill, draw=figAdapt, minimum width=2.7cm, anchor=east]
      (cli) at (-5.15,-0.92) {Consola};
    \node[figbox, fill=figAdaptFill, draw=figAdapt, minimum width=2.7cm, anchor=west]
      (sql) at (5.15,0.92) {SQL Server};
    \node[figbox, fill=figAdaptFill, draw=figAdapt, minimum width=2.7cm, anchor=west]
      (mongo) at (5.15,-0.92) {MongoDB};
    \draw[figarr] (api.east) -- (pin.north west);
    \draw[figarr] (cli.east) -- (pin.south west);
    \draw[figarr] (sql.west) -- (pout.north east);
    \draw[figarr] (mongo.west) -- (pout.south east);
    \node[font=\scriptsize, text=figMuted, anchor=north] at (0,-2.72)
      {Las flechas apuntan al núcleo: el adaptador depende del puerto, no al revés.};
  \end{tikzpicture}
  \caption{Puertos y adaptadores. La API y la consola usan el mismo puerto de entrada; SQL Server y MongoDB implementan el de salida. El núcleo no distingue el canal ni el motor. Elaboración propia a partir de \parencite{cockburn2005hexagonal}.}
  \label{fig:hexagonal}
\end{figure}

\subsubsection{¿Cuál sería el impacto en la testabilidad y la capacidad de cambiar la base de datos de SQL Server a MongoDB sin afectar la lógica de negocio?}

Con dependencias hacia adentro, los casos de uso y las entidades se prueban con \textbf{dobles de
prueba} en memoria (repositorios falsos, relojes controlados, buses de eventos simulados) sin levantar
base de datos ni servidor web \parencite{martin2017clean,cockburn2005hexagonal}. Las pruebas dejan
de ser de integración pesada por defecto y pasan a ser unitarias rápidas y deterministas.

Además, al no filtrar detalles de persistencia al dominio, los tests dejan de romperse por cambios
cosméticos en columnas o en el framework de UI: solo fallan cuando se altera un contrato de puerto
o una regla de negocio explícita (figura~\ref{fig:testabilidad}).

\begin{figure}[H]
  \centering
  \begin{tikzpicture}[node distance=6.5mm]
    \node[figbox, fill=figBadFill, draw=figBad, minimum width=5.15cm] (a1) {Prueba};
    \node[font=\footnotesize\bfseries, text=figBad, anchor=south] (ha)
      at ([yshift=1.6mm]a1.north) {Acoplada};
    \node[figbox, fill=figBadFill, draw=figBad, minimum width=5.15cm, below=of a1] (a2) {Servicio de inventario};
    \node[figbox, fill=figInfraFill, draw=figInfra, minimum width=5.15cm, below=of a2] (a3) {%
      SQL Server\\y servidor web};
    \node[font=\scriptsize\bfseries, text=figBad, below=1.5mm of a3] (afoot) {lenta y frágil};
    \draw[figarrbad] (a1) -- (a2);
    \draw[figarrbad] (a2) -- (a3);

    \node[figbox, fill=figOkFill, draw=figOk, minimum width=5.35cm, right=12mm of a1] (b1) {Prueba};
    \node[font=\footnotesize\bfseries, text=figOk, anchor=south] (hb)
      at ([yshift=1.6mm]b1.north) {Con puertos};
    \node[figbox, fill=figNucleoFill, draw=figNucleo, minimum width=5.35cm, below=of b1] (b2) {%
      \texttt{\scriptsize RegistrarEntradaDeStock}};
    \node[figbox, fill=figOkFill, draw=figOk, minimum width=5.35cm, below=of b2] (b3) {%
      Doble en memoria\\
      {\scriptsize repositorio, reloj, bus de eventos}};
    \node[font=\scriptsize\bfseries, text=figOk, below=1.5mm of b3] (bfoot) {rápida y determinista};
    \draw[figarrook] (b1) -- (b2);
    \draw[figarrook] (b2) -- (b3);
  \end{tikzpicture}
  \caption{Efecto en la testabilidad. A la izquierda la prueba arrastra la base de datos y el servidor web. A la derecha el caso de uso se ejerce con dobles en memoria y no se entera de un cambio cosmético de columnas. }
  \label{fig:testabilidad}
\end{figure}

Supongamos que el dominio modela \emph{Producto}, \emph{Almacén} y \emph{Movimiento} con invariantes
(propiedades de stock no negativo, reservas coherentes, etc.). Bajo N-Capas acopladas, migrar a
MongoDB implica retocar servicios, mapeos, validaciones y posiblemente pantallas que asumían
identificadores o joins relacionales.

Con Regla de Dependencia y puertos/adaptadores:

\begin{enumerate}
  \item Se define (o se mantiene) el puerto \texttt{IRepositorioDeInventario} en el núcleo.
  \item El adaptador \texttt{SqlServerInventarioAdapter} implementa el puerto con tablas y
    transacciones ACID.
  \item Un nuevo \texttt{MongoInventarioAdapter} implementa el mismo puerto con colecciones y,
    si hace falta, transacciones multi-documento; el mapeo documento--entidad vive solo en el
    adaptador.
  \item La composición (inyección de dependencias en el arranque de la aplicación) selecciona el
    adaptador según configuración, sin recompilar el dominio.
\end{enumerate}

La lógica de negocio permanece idéntica; el coste de migración se acota a infraestructura, pruebas
de adaptador y estrategia de datos (dual write, backfill, corte), no a reescribir reglas de
inventario (figura~\ref{fig:migracion}) \parencite{richards2015patterns,martin2017clean}.

\begin{figure}[H]
  \centering
  \begin{tikzpicture}
    \node[draw=figOk, fill=figOkFill, rounded corners=2pt, line width=0.8pt,
          minimum width=6.45cm, minimum height=3.85cm] (stay) {};
    \node[anchor=north, font=\footnotesize\bfseries, text=figOk]
      at ([yshift=-2.4mm]stay.north) {No se recompila};
    \node[anchor=north west, font=\scriptsize, align=left]
      at ([xshift=3.2mm,yshift=-8.2mm]stay.north west)
      {Entidades: Producto, Almacén\\
       y Movimiento.\\[3pt]
       Caso de uso:\\
       \texttt{RegistrarEntradaDeStock}.\\[3pt]
       Puerto:\\
       \texttt{IRepositorioDeInventario}.};

    \node[draw=figInfra, fill=figInfraFill, rounded corners=2pt, line width=0.8pt,
          minimum width=6.45cm, minimum height=3.85cm, right=5mm of stay] (swap) {};
    \node[anchor=north, font=\footnotesize\bfseries, text=figInfra]
      at ([yshift=-2.4mm]swap.north) {Se sustituye};
    \node[anchor=north west, font=\scriptsize, align=left]
      at ([xshift=3.2mm,yshift=-8.2mm]swap.north west)
      {\texttt{SqlServerInventarioAdapter}\\
       sale del sistema.\\[3pt]
       Entra \texttt{MongoInventarioAdapter}.\\[3pt]
       Colecciones, mapeo y\\
       transacciones del motor.\\[3pt]
       Se rehacen las pruebas del adaptador,\\
       no las reglas.};

    \coordinate (midpanels) at ($(stay.south)!0.5!(swap.south)$);
    \node[figbox, fill=figUiFill, minimum width=13.4cm, anchor=north]
      (boot) at ([yshift=-5.5mm]midpanels) {%
      Arranque: la composición elige el adaptador por configuración};
    \draw[figarr] (boot.north -| stay.center) -- (stay.south);
    \draw[figarr] (boot.north -| swap.center) -- (swap.south);
  \end{tikzpicture}
  \caption{Límite del cambio de SQL Server a MongoDB. El núcleo y el puerto permanecen; el coste queda en el adaptador, en sus pruebas y en la estrategia de datos. }
  \label{fig:migracion}
\end{figure}

\subsubsection{Síntesis}

El error no es haber elegido N-Capas al inicio, sino permitir que las dependencias del dominio
apunten hacia afuera hacia detalles volátiles. La Regla de Dependencia y los puertos/adaptadores
no añaden capas por capricho: \textbf{invierten el flujo de control} para que el negocio mande sobre
la tecnología. El resultado es un sistema más costoso de diseñar al principio, pero con pruebas
rápidas y con la libertad de sustituir SQL Server por MongoDB (u otro almacén) sin erosionar el
núcleo del inventario.
